\documentclass{tietreport}

% Import images
\usepackage{graphicx}

% Bibliography configuration
\usepackage[
  date=year,
  style=alphabetic,
  backend=biber,
  sorting=nty,
  maxnames=5,
  minnames=3,
  maxalphanames=4,
  minalphanames=3,
  backref=true,
  doi=false,
  isbn=false,
  url=false,
  eprint=false
]{biblatex}
\addbibresource{references.bib}

%% ==================================================
%% CUSTOMIZE THESE FIELDS FOR YOUR PROJECT
%% ==================================================

% Course/Lab header
\TitlePageHeader{%
  UCS503 - Software Engineering%
}

% Main project title
\ProjectTitle{%
  EduBridge%
}

% Subtitle or project description
\ProjectSubTitle{%
  A Progressive Web Application for Small-Scale Coaching Institute Management%
}

% Student information (up to 4 students)
\author{%
  \texttt{1024170013} Abhijay Nagal \\
  \texttt{1024170120} Vansh Wadhwa \\
  \texttt{1024170115} Samarth Gulati \\
  \texttt{1024170010} Brahmansh Singh Sahni%
}

% Group number, degree, year, program
\TitlePageSubText{%
  Group: \texttt{3Q14} BE Third Year, CSE%
}

% Faculty advisor/guide name
\AdvisorName{%
  Sandeep Kaur%
}

% Evaluation period and department info
\TitlePageFooterText{{%
    \large Mid-Semester Evaluation \\[0.2cm]
    \bfseries Computer Science and Engineering
    Department%
  } \\[0.15cm]
  Thapar Institute of Engineering and Technology,
  Patiala%
}

%% ==================================================

\begin{document}

% Generate title page
\maketitle

% Table of contents (auto-generated from chapters)
\tableofcontents

% ===== CHAPTER 1: INTRODUCTION =====

\chapter{Introduction}

\section{Background and Context}

Small tuition centres and standalone coaching institutes make up a large part
of supplementary education in India, but almost none of the learning
management systems on the market are built for them. Products such as Google
Classroom or Moodle assume a school-sized IT setup and administrative staff to
run it. A tuition centre with one or two teachers and a few dozen students
does not have either. In practice, these institutes run on WhatsApp groups for
announcements, phone calls for attendance and fee reminders, and paper
registers for records. This works, but it does not scale past a certain
number of students and it leaves parents with almost no visibility into how
their child is actually doing.

\subsection{Problem Statement}

The specific problem this project addresses is the absence of a lightweight,
role-appropriate platform for three distinct groups of users at a small
coaching institute: students, who need access to study material,
quizzes and their own attendance and academic record; parents, who need to
see their child's attendance, progress and fee status without having to call
the teacher; and administration, who need to run the entire operation, upload
material, mark attendance, record marks, track fees and broadcast
notifications, from a single screen. Existing communication over WhatsApp is
unstructured, not searchable, and puts the entire burden of manually sending
reminders on the teacher. EduBridge replaces this with a single application
that keeps this data structured and makes it available to the right person
automatically.

\section{Project Objectives}

\begin{enumerate}
  \item \textbf{Objective 1:} Design and build a single web-based platform
    with three role-based logins, student, parent and administration, that
    lets a small coaching institute manage attendance, study material,
    quizzes, marks and fee records without relying on WhatsApp groups or
    phone calls, delivered as a working prototype within the semester.
  \item \textbf{Objective 2:} Implement the core administrative workflows,
    batch and student management, attendance marking, quiz creation with
    automatic grading, marks entry and fee tracking, so that any data the
    administration enters is immediately reflected on the student and parent
    dashboards without a manual sync step.
  \item \textbf{Objective 3:} Package and deploy the application as an
    installable progressive web app hosted on Vercel with continuous
    deployment from GitHub, so that every push to the main branch is
    automatically built and released without a manual deployment step.
\end{enumerate}

% ===== CHAPTER 2: METHODOLOGY & DESIGN =====

\chapter{Methodology and Design}

\section{System Architecture}

\subsection{High-Level Design}

EduBridge is a single-page application with one shared authentication system
rather than three separate logins. Every account sits in one users table with
a role field (\texttt{student}, \texttt{parent} or \texttt{admin}), and the
application routes each role to its own dashboard and navigation after login.
This keeps the authentication logic in one place while still giving each role
a completely separate experience once inside the app.

All screens read and write through a single data-access module rather than
talking to storage directly. This module currently persists data in the
browser using IndexedDB, seeded with a realistic demo dataset on first launch,
which means the whole prototype runs with no backend server and can be
evaluated by opening the deployed link with no setup. Because every screen
goes through this one module, moving to a shared backend at a later stage
means rewriting that module only, not the components built on top of it.

Figure~\ref{fig:usecase} shows the use case diagram for the system, covering
the three actors and the behaviour available to each of them, including the
include and extend relationships between login, quiz evaluation and the
notification and fee reminder flows. Figure~\ref{fig:activity} shows the
activity diagram for the login and role-dispatch flow, including the
credential validation loop, the role-based split into the three dashboards,
and the logout decision.

\begin{figure}[h]
  \centering
  \includegraphics[width=0.95\textwidth]{figures/usecase_diagram.jpg}
  \caption{Use case diagram for EduBridge, showing the Student, Parent and
    Administration actors and their associated use cases.}
  \label{fig:usecase}
\end{figure}

\begin{figure}[h]
  \centering
  \includegraphics[width=0.8\textwidth]{figures/activity_diagram.jpg}
  \caption{Activity diagram for the login and role-dispatch flow, showing the
    credential validation, role-based dashboard split and logout decision.}
  \label{fig:activity}
\end{figure}

\subsection{Technical Stack}

\begin{itemize}
  \item \textbf{Framework/Language:} React~\cite{react2024} built with
    Vite~\cite{vite2024}, written in plain JavaScript. Routing is handled by
    React Router.
  \item \textbf{Database:} IndexedDB, accessed through the
    \texttt{idb-keyval}~\cite{idbkeyval} helper library, behind a single
    data-access module.
  \item \textbf{Tools:} Git and GitHub for version control, Vercel for
    hosting with continuous deployment, \texttt{sharp} for generating PWA
    icons at build time.
  \item \textbf{Libraries:} Tailwind CSS~\cite{tailwind2024} for styling,
    Recharts for the attendance and progress charts,
    \texttt{vite-plugin-pwa} for the service worker and installability,
    following the Progressive Web App standard~\cite{mdn_pwa}.
\end{itemize}

\section{Implementation Details}

\subsection{Module 1: Core Functionality}

The authentication module checks a phone number and password against the
users table and stores the resulting session in local storage, with route
guards that redirect a signed-in user to their own role's dashboard and log
out any user whose underlying student or parent record has been removed.

The attendance module lets the administration mark an entire batch in one
screen, with a three-state toggle per student (present, late, absent) and an
"all present" shortcut, since a teacher marking thirty students by hand needs
this to be fast. Saving overwrites any existing record for that batch and
date rather than creating duplicates, which was verified directly with an
automated test.

The quiz module supports multiple-choice questions with a single correct
option per question. A student can attempt a quiz once; on submission the
module computes the score against the stored correct answers and stores the
attempt, after which the student sees a full answer review showing which
options were correct and which they chose.

\subsection{Module 2: Secondary Features}

The notification module lets administration post an announcement targeted at
everyone, students only or parents only, and optionally scoped to a single
batch. Each user's notification list is filtered by their role and batch
membership, so a parent never sees a student-only notice and vice versa. An
unread count is tracked per user and shown as a badge on the navigation bar.

The fee module tracks a due amount and date per student, automatically flags
an entry as overdue once its due date passes, and lets administration record
a payment against it. A single button posts one reminder notification to all
parents with an overdue fee, which is the manual version of the scheduled
reminder job described in the original project plan.

The report card module stores marks and a remark per subject per test, and
renders the same data as a trend line across tests and a bar chart per
subject on both the student's own progress screen and the parent's view of
their child, using Recharts.

Because the whole application is registered as a service worker with a web
app manifest, it can be installed to a phone's home screen from the browser
and opens without any browser chrome, which was a specific requirement for
this project.

% ===== CHAPTER 3: RESULTS & EVALUATION =====

\chapter{Results and Evaluation}

\section{Testing Methodology}

\begin{itemize}
  \item \textbf{Unit Testing:} The data-access module was exercised outside
    the browser with a Node script that stubs IndexedDB in memory. It covers
    authentication with correct and incorrect credentials, role-based
    notification filtering, quiz scoring on both a fully correct and a fully
    incorrect attempt, attendance overwrite behaviour, fee overdue flagging,
    payment recording, and cascading deletes when a student or batch is
    removed.
  \item \textbf{Integration Testing:} The production build was served
    locally and every top-level route, including deep links such as
    \texttt{/student/quizzes}, was checked to confirm it returns a valid page
    rather than a routing error, which specifically verifies that the
    single-page application rewrite rule on Vercel is configured correctly.
  \item \textbf{Performance Testing:} The production bundle was inspected
    after enabling manual code-splitting, separating the charting library
    into its own chunk so it is not loaded on screens that do not need it.
  \item \textbf{User Testing:} Not carried out yet at this stage of the
    project. This is listed under future work.
\end{itemize}

\section{Performance Metrics}

Table~\ref{tab:performance} reports figures taken directly from the
production build and from the automated checks described above, rather than
simulated targets.

\begin{table}[h]
\centering
\begin{tabular}{|l|r|}
  \hline
  \textbf{Metric} & \textbf{Result} \\
  \hline
  Automated data-layer checks passed & 38 / 38 \\
  \hline
  Core routes returning a valid page & 7 / 7 \\
  \hline
  Main application bundle (gzipped) & 25.9 kB \\
  \hline
  Charting library bundle (gzipped, lazy-loaded) & 105.6 kB \\
  \hline
  React and routing bundle (gzipped) & 53.5 kB \\
  \hline
  PWA icon sizes generated & 192px, 512px, maskable 512px \\
  \hline
\end{tabular}
\caption{Build and test results for the current prototype}
\label{tab:performance}
\end{table}

\section{Analysis}

The first two objectives, a working three-role platform and the core
administrative workflows, have been met in full for this evaluation. The
third objective, packaging and deployment, was also met and slightly
exceeded: the application is live on Vercel with continuous deployment
already wired up, and it is installable as a home screen app on both Android
and iOS, which was treated as a stretch goal at the planning stage rather
than a requirement.

The main challenge encountered was choosing a data layer for this stage of
the project. The original plan called for a shared Supabase backend, but
setting that up correctly, including authentication rules and row-level
security, was judged too large a task to complete safely in the time
available before this evaluation without risking a broken submission. The
decision taken was to build a fully functional local data layer instead,
isolated behind one module, so that every feature works end to end for
grading and demonstration purposes, with the understanding that moving to a
shared backend is a contained follow-up task rather than a rewrite. The
direct consequence of this decision is that data is currently local to each
browser rather than shared across devices, which is documented as a known
limitation rather than left as a silent gap.

% ===== CHAPTER 4: CONCLUSION =====

\chapter{Conclusion}

\section{Summary of Work}

This evaluation covers the full first version of EduBridge: a three-role
progressive web application covering study material, quizzes, attendance,
report cards, fee tracking and targeted notifications, deployed with
continuous integration from GitHub to Vercel. All three stated objectives
were met within the current stage of the project.

\section{Key Findings}

\begin{itemize}
  \item \textbf{Finding 1:} Modelling all three roles around one
    authentication table with a role field, rather than three separate login
    systems, made the role-based routing and access rules considerably
    simpler to implement and to reason about than three independent systems
    would have been.
  \item \textbf{Finding 2:} Routing every screen through a single data-access
    module, rather than letting individual components read storage directly,
    made it possible to write meaningful automated tests for the application
    logic even without a browser test environment available in this setup.
  \item \textbf{Finding 3:} A local-first data layer was sufficient to
    deliver a fully working, installable prototype without any backend
    infrastructure, which was not the original plan but turned out to be a
    reasonable way to de-risk the delivery timeline for this evaluation.
\end{itemize}

\section{Future Work and Recommendations}

\begin{enumerate}
  \item Replace the local IndexedDB data layer with a shared Supabase
    backend for authentication, the database and file storage, so that data
    is consistent across every device rather than local to one browser.
  \item Add real file uploads for study material, since the current version
    only supports written notes and external links.
  \item Integrate an SMS or WhatsApp Business API channel for fee reminders
    and class-cancellation notices, since in-app notifications alone are
    unlikely to reach parents who do not open the app daily.
  \item Extend automated testing to cover the user interface itself once a
    browser testing setup is available, rather than only the underlying data
    logic.
  \item Carry out a short round of user testing with an actual tuition
    teacher and a small batch of students to validate the attendance marking
    and quiz flows under real conditions before any real deployment.
\end{enumerate}

\section{Final Remarks}

Building EduBridge around a single, swappable data-access layer turned out to
be the most useful early decision in the project. It meant the team could
commit to a local-first version for this stage without locking the
architecture into that choice permanently, and it kept every screen simple to
build because none of them needed to know how or where the data was actually
stored.

% ===== BIBLIOGRAPHY =====

% Print bibliography with heading in TOC
\printbibliography[heading=bibintoc]

\end{document}
